\documentclass[12pt, a4paper]{article}

\usepackage[utf8]{inputenc}
\usepackage[T1]{fontenc}
\usepackage[spanish]{babel}
\usepackage{amsmath}
\usepackage{amssymb}
\usepackage{geometry}
\usepackage{listings}
\usepackage{xcolor}
\usepackage{parskip}

\geometry{margin=2.5cm}

\lstdefinestyle{avr}{
    language=C,
    basicstyle=\ttfamily\small,
    keywordstyle=\color{blue}\bfseries,
    commentstyle=\color{gray}\itshape,
    stringstyle=\color{teal},
    numbers=left,
    numberstyle=\tiny\color{gray},
    frame=single,
    breaklines=true,
    tabsize=4,
    showstringspaces=false,
}

\title{Ejercicio 3 --- Análisis Teórico\\
\large Trabajo Práctico 1: Interrupciones y ADC por Polling}
\author{Sistemas de Tiempo Real}
\date{}

\begin{document}

\maketitle

% =====================================================================
\section{Parámetros Temporales y Frecuencias}
% =====================================================================

\subsection{Prescaler del ADC}

El datasheet del ATmega328P especifica que el reloj del ADC ($F_{ADC}$) debe
estar comprendido entre 50~kHz y 200~kHz para garantizar los 10~bits de
resolución. Con un cristal de 16~MHz ($F_{CPU}$) se selecciona un prescaler de
128 (\texttt{ADPS2:0 = 111}):

\[
    F_{ADC} = \frac{F_{CPU}}{\text{Prescaler}} = \frac{16\,000\,000}{128} = 125\text{ kHz}
\]

El valor de 125~kHz se ubica en el centro del rango admisible, maximizando el
margen de tolerancia frente a variaciones de tensión o temperatura.

\subsection{Tiempo de Conversión}

Una conversión normal (posterior a la primera) toma exactamente 13 ciclos de
reloj del ADC:

\[
    T_{conv} = 13 \cdot \frac{1}{F_{ADC}} = 13 \cdot \frac{1}{125\,000} = 104\;\mu\text{s}
\]

Durante este intervalo la CPU permanece en espera activa (\emph{polling}),
bloqueada hasta que el ADC levante la bandera \texttt{ADIF}.

\subsection{Resolución Analógica}

Con referencia de tensión $V_{REF} = 5$~V y convertidor de $n = 10$~bits:

\[
    \text{Resolución} = \frac{V_{REF}}{2^{n}} = \frac{5\text{ V}}{1024} \approx 4{,}88\text{ mV por bit}
\]

% =====================================================================
\section{Análisis del Código}
% =====================================================================

\subsection{Polling de bandera (\texttt{adc\_read})}

El enunciado requiere explícitamente \emph{consulta de bandera (polling)}. La
implementación utiliza espera activa:

\begin{lstlisting}[style=avr]
while (!(ADCSRA & (1 << ADIF)));
ADCSRA |= (1 << ADIF);
\end{lstlisting}

La espera activa es aceptable en este contexto porque el sistema es
\emph{bare-metal} de lazo único: no existe otra tarea bloqueada durante los
104~µs de conversión. El comportamiento es completamente determinista.

La bandera \texttt{ADIF} sigue la convención \textbf{Clear-on-Write-1} de la
arquitectura AVR: se limpia escribiendo un \texttt{1} lógico sobre el bit.

\subsection{Atomicidad en la lectura (\texttt{adc\_init} y lectura)}

Al configurar \texttt{ADLAR = 0} mediante:

\begin{lstlisting}[style=avr]
ADMUX = (1 << REFS0);
\end{lstlisting}

los 10 bits de resultado quedan alineados a la derecha en los registros
\texttt{ADCL}/\texttt{ADCH}. La instrucción \texttt{return ADC;} garantiza
atomicidad: el compilador AVR-GCC lee primero \texttt{ADCL} (lo que bloquea
por hardware la actualización del par de registros) y luego \texttt{ADCH},
asegurando la integridad del dato de 16~bits.

% =====================================================================
\section{Multiplexado y Tiempos de Latching}
% =====================================================================

\subsection{Manipulación a nivel de bits}

Para separar los nibbles destinados a cada latch 74LS373:

\begin{lstlisting}[style=avr]
uint8_t byte_alto = (uint8_t)((valor >> 8) & 0x0F);
uint8_t byte_bajo = (uint8_t)(valor & 0xFF);
\end{lstlisting}

\begin{itemize}
    \item \textbf{Desplazamiento (\texttt{>> 8}):} mueve los dos bits más
          significativos del resultado de 10~bits a las posiciones 0 y~1.
    \item \textbf{Enmascaramiento (\texttt{\& 0x0F}):} fuerza a cero los cuatro
          bits altos del byte, evitando activar pines no deseados del
          \texttt{PORTB} conectados a hardware externo.
\end{itemize}

\subsection{Tiempos de latch (Setup y Hold Time)}

\begin{lstlisting}[style=avr]
PORTC |= (1 << PORTC2);   // LE alto
_delay_us(5);
PORTC &= ~(1 << PORTC2);  // LE bajo
\end{lstlisting}

El retardo de 5~µs no es arbitrario: cumple el \emph{Pulse Width} mínimo y el
\emph{Data Setup Time} especificados en el datasheet del 74LS373. Si el
microcontrolador conmutara la señal \texttt{LE} a la frecuencia del sistema
($\approx 62{,}5$~ns por ciclo), las compuertas TTL no alcanzarían a
estabilizarse, dando lugar a capturas erróneas por capacitancia parásita.

% =====================================================================
\section{Período de Muestreo}
% =====================================================================

El lazo principal incluye \texttt{\_delay\_ms(50)}. El período real de cada
iteración es la suma de todas las operaciones:

\[
    T_{ciclo} = T_{conv}(104\;\mu\text{s})
              + T_{latch\,alto}(5\;\mu\text{s})
              + T_{latch\,bajo}(5\;\mu\text{s})
              + 50\text{ ms}
              + t_{overhead}
\]

El retardo de 50~ms fija una tasa de refresco de 20~Hz, suficiente para la
persistencia visual del ojo humano sin generar \emph{aliasing} perceptible en
el display de 7~segmentos.

La implementación mediante \texttt{while(1)} con retardo bloqueante introduce
\emph{jitter} inherente. Para un sistema de tiempo real estricto se debería
emplear un temporizador por interrupción; sin embargo, el enunciado requiere
explícitamente el modo de muestreo por bucle y consulta de bandera, y la
tolerancia de una interfaz humana acepta dicha fluctuación.

\end{document}